\itemhead{4.1 例题 · 例题9（教材p.102）}{XJ-B1-C04-EX-P102-Y0669}
\label{XJ-B1-C04-EX-P102-Y0669}
{\footnotesize 湘教版教材p.102 / PDF第109页 / 4.1 例题 / 例题9}
\field{原题}{已知幂函数 $f(x)$ 满足 $f(8)=16$，求 $f(9)$ 和 $f(64)$ 的值。}
\field{参考答案}{设 $f(x)=x^\alpha$，$8^\alpha=16$ 得 $\alpha=4/3$；$f(9)=9^{4/3}=9\sqrt[3]9$，$f(64)=256$。}
{\small \field{关联知识}{P03 幂函数的概念、图象与性质（教材帮p.125、126、127 / PDF134、135、136）}
\field{方法与易错点}{设 $f(x)=x^\alpha$，用已知点求指数，再查该指数的定义域、奇偶性和各单调区间。 易错：把 $3x^2$ 当标准幂函数；将第一象限性质直接用于负半轴；分数负指数遗漏0禁点。}
\field{答案依据}{例题原解位置：湘教版教材p.102及题库原图；独立编写可编辑参考解答。答案册此部分主要收练习、习题，不伪造例题答案册位置。}}
